\subsection{全息/塌缩、Mellin 主核、奇偶补账与冻结层析的分层闭包}\label{subsec:conclusion-holography-collapse-mellin-parity-freezing-layered-closure}
沿用定理
\ref{thm:conclusion-fiber-toggle-boolean-spectrum-polynomial}、
\ref{thm:conclusion-fiber-toggle-sector-polynomial-complete-reconstruction}、
\ref{cor:conclusion-fiber-toggle-sector-polynomial-recovers-geometry-symmetry}、
\ref{thm:conclusion-fiber-toggle-boolean-function-algebra-collapse}、
\ref{thm:conclusion-fiber-fixed-diameter-multiplicity-extremals}、
\ref{thm:conclusion-fiber-fixed-diameter-symmetry-extremals}、
\ref{thm:conclusion-fiber-fixed-diameter-phi-baseline-gap}、
\ref{thm:conclusion-full-residue-regular-simplex-orbit}、
\ref{cor:conclusion-full-residue-phi-square-universal-exponent}、
\ref{thm:conclusion-foldbin-highgenus-minsector-tomography}、
\ref{cor:conclusion-foldbin-audited-even-fibonacci-minsector-genus-recovery}、
\ref{thm:conclusion-frozen-escort-shannon-renyi-collapse}、
\ref{thm:conclusion-freezing-infinity-budget-noncommutation}、
\ref{thm:conclusion-sample-pattern-family-closes-hot-contact-geometry}、
\ref{thm:conclusion-capacity-curve-order-holography-universality}、
\ref{thm:conclusion-majorization-capacity-infonce-strict-antiisomorphism}、
\ref{thm:conclusion-mellin-master-kernel-two-sided-completeness}、
\ref{thm:conclusion-thermodynamic-flattening-subextensive-arithmetic-rigidity}、
\ref{thm:conclusion-oddwindow-full-inversion-evenwindow-rankone-parity-blindmode}、
\ref{thm:conclusion-first-coordinate-bias-global-collision-resonance-homology}
与定理
\ref{thm:conclusion-finite-defect-profile-minimal-cauchy-rank}
的记号。前述结论已经分别把纤维扇区维数多项式对长度谱与几何/对称的全恢复、固定直径层上的链相/超立方相双端极值、原子主导 residue-law 的单纯形轨道、冻结极限与容量端点的不交换性、样本模式对热相接触几何的闭合、连续容量曲线与完整 Bayes--InfoNCE 塔的互逆完备性、Mellin--Dirichlet 主核的双侧统一、广延热力学与 even-\(\zeta\) 算术的二层分裂、奇/偶窗口上的局部可逆性分界，以及有限缺陷剖面的最小 Cauchy 秩固定下来。把这些对象再作一次跨节闭包，可抽出如下十条终端结论。

\begin{theorem}[扇区维数多项式与抽象开关不变代数之间的全息/塌缩二层分裂]\label{thm:conclusion-sector-polynomial-vs-toggle-algebra-holography-collapse-splitting}
\leanverified{paper\_conclusion\_sector\_polynomial\_vs\_toggle\_algebra\_holography\_collapse\_splitting}
对任意稳定类型 \(x\)，记
\[
P_x(t):=\prod_{i=1}^{r(x)}\bigl(1+(n_i-1)t\bigr),
\qquad
n_i=F_{\ell_i(x)+2}.
\]
则 \(P_x(t)\) 唯一决定长度谱
\[
\mathrm{Spec}^\ast(x):=\{\ell_i(x)\}_{i=1}^{r(x)},
\]
从而唯一决定
\[
\mathcal F_x,\qquad
\mathcal T_x,\qquad
\Aut(\mathcal F_x),\qquad
|\mathcal F_x|,\qquad
\mathrm{diam}(\mathcal F_x),\qquad
\dim_\square(\mathcal F_x).
\]
与此相反，抽象开关不变代数满足
\[
\End(\mathcal H_x)^{\mathcal T_x}
\cong
\mathrm{Fun}(\{0,1\}^{r(x)},\CC)
\cong
\CC^{2^{r(x)}},
\]
其抽象同构型只依赖于 Cartesian 素因子个数 \(r(x)\)，与 \(\Spec^\ast(x)\) 的更细分拆无关。

因此，存在一条严格的忘却态射
\[
P_x
\longmapsto
\End(\mathcal H_x)^{\mathcal T_x},
\]
它保留全部布尔扇区数，却抹去全部长度数据；换言之，\(P_x\) 是完整几何--对称全息坐标，而 \(\End(\mathcal H_x)^{\mathcal T_x}\) 是其终端布尔塌缩商。
\end{theorem}
\begin{proof}
扇区维数多项式的闭式与其对 \(\{n_i\}\)、\(\mathrm{Spec}^\ast(x)\) 的唯一反演，正是定理 \ref{thm:conclusion-fiber-toggle-boolean-spectrum-polynomial} 与定理 \ref{thm:conclusion-fiber-toggle-sector-polynomial-complete-reconstruction} 的内容。由推论 \ref{cor:conclusion-fiber-toggle-sector-polynomial-recovers-geometry-symmetry}，一旦 \(\mathrm{Spec}^\ast(x)\) 固定，\(\mathcal F_x\)、\(\mathcal T_x\)、\(\Aut(\mathcal F_x)\)、\(|\mathcal F_x|\)、\(\mathrm{diam}(\mathcal F_x)\) 与 \(\dim_\square(\mathcal F_x)\) 便全部唯一恢复。

另一方面，定理 \ref{thm:conclusion-fiber-toggle-boolean-function-algebra-collapse} 已证明
\[
\End(\mathcal H_x)^{\mathcal T_x}\cong \CC^{2^{r(x)}},
\]
并且该代数的抽象同构型当且仅当 \(r(x)\) 相同。故它只记住布尔扇区数，而不记住长度分拆本身。证毕。
\end{proof}

\begin{theorem}[固定直径层中链相与超立方相构成乘法型审计不变量的共同双端相]\label{thm:conclusion-fixed-diameter-chain-cube-common-dual-phase}
\leanverified{paper\_conclusion\_fixed\_diameter\_chain\_cube\_common\_dual\_phase}
固定直径
\[
\mathrm{diam}(\mathcal F_x)=D\ge 1.
\]
则同时有
\[
F_{D+2}\le d_m(x)\le 2^D,
\qquad
2\le |\Aut(\mathcal F_x)|\le 2^D D!,
\]
并且对
\[
\Delta_x:=\log d_m(x)-D\log\varphi
\]
有
\[
\log F_{D+2}-D\log\varphi
\le
\Delta_x
\le
D\log(2/\varphi).
\]
三组上下界分别且唯一地由
\[
\Gamma_D
\qquad\text{与}\qquad
Q_D=(\Gamma_1)^D
\]
取到：链相 \(\Gamma_D\) 同时最小化多重度、对称体积与 \(\varphi\)-基线偏离；超立方相 \(Q_D\) 同时最大化这三者。

特别地，
\[
\log F_{D+2}-D\log\varphi \to \log(\varphi^2/\sqrt5)
\]
始终有界，而
\[
D\log(2/\varphi)
\]
则按 \(D\) 线性增长。故固定直径层不是对称相图，而是一条由链相与超立方相支配的单边线性放大相图。
\end{theorem}
\begin{proof}
多重度的双端极值由定理 \ref{thm:conclusion-fiber-fixed-diameter-multiplicity-extremals} 给出，而对称群阶的双端极值由定理 \ref{thm:conclusion-fiber-fixed-diameter-symmetry-extremals} 给出。定理 \ref{thm:conclusion-fiber-fixed-diameter-phi-baseline-gap} 则进一步证明：\(\Delta_x\) 的上下界也恰由同一对极值对象 \(\Gamma_D\) 与 \(Q_D\) 唯一达到，并给出下端有界、上端线性增长的精确渐近。三者合并即得结论。证毕。
\end{proof}

\begin{theorem}[原子主导 residue-law 的角向正则单纯形与径向 \texorpdfstring{$\varphi^{-2}$}{phi^-2} 冻结]\label{thm:conclusion-atomic-residue-simplex-radial-freeze-splitting}
若处于原子主导区
\[
\rho_m^{\mathrm{core}}<1,
\]
定义偶长度归一化 residue-law 偏差
\[
\nu_q:=\mu_{m,\bullet}^{\mathrm{full}}(2q)-u_m,
\qquad
u_m=\frac1m\mathbf 1.
\]
则存在顶点集
\[
V_m=\{v_r:r\in\ZZ/m\ZZ\}\subset
H_m:=
\Bigl\{x\in\RR^{\ZZ/m\ZZ}:\ \sum_r x_r=0\Bigr\}
\]
使得
\[
\frac{a_{2q}^{\mathrm{full}}(1)}{2}\,\nu_q
=
v_{q\bmod m}+o(1),
\]
且 \(V_m\) 构成 \(H_m\) 中的正则 \(A_{m-1}\)-单纯形：
\[
\langle v_r,v_s\rangle=\delta_{rs}-\frac1m,
\qquad
\|v_r\|_2^2=\frac{m-1}{m},
\qquad
\|v_r-v_s\|_2^2=2\ \ (r\neq s).
\]
并且经验测度收敛为
\[
\frac1N\sum_{q=0}^{N-1}\delta_{\frac{a_{2q}^{\mathrm{full}}(1)}{2}\nu_q}
\Rightarrow
\frac1m\sum_{r\in\ZZ/m\ZZ}\delta_{v_r}.
\]
与此同时，奇长度偏差经归一化后塌缩为零，而整个 residue-law 的统一指数满足
\[
\limsup_{n\to\infty}
\|\mu_{m,\bullet}^{\mathrm{full}}(n)-u_m\|_2^{1/n}
=
\varphi^{-2},
\]
该指数与模数 \(m\) 无关。

因此，原子主导区中的 residue-law 不是一般零均值云，而是“模数决定角向正则单纯形、\(\varphi^{-2}\) 决定径向冻结”的刚性双层对象。
\end{theorem}
\begin{proof}
正则单纯形轨道与经验测度极限由定理 \ref{thm:conclusion-full-residue-regular-simplex-orbit} 直接给出。奇长度塌缩与统一指数冻结则正是推论 \ref{cor:conclusion-full-residue-phi-square-universal-exponent} 的内容。证毕。
\end{proof}

\begin{theorem}[高 genus 极限与零温冻结分别层析最小扇区与最大扇区，而 \texorpdfstring{$\infty$}{infty} 预算端点不交换]\label{thm:conclusion-highgenus-freezing-infinity-budget-dual-tomography-split}
设
\[
d_{\min}(m):=\min_x d_m(x),
\qquad
N_{\min}(m):=\#\{x:d_m(x)=d_{\min}(m)\},
\]
以及最大纤维集
\[
A_m^{\max}:=\{x:d_m(x)=M_m\}.
\]
则高 genus 极限满足
\[
d_{\min}(m)^{2g-2}Z_m(\Sigma_g)\longrightarrow N_{\min}(m)
\qquad(g\to\infty),
\]
并且在审计偶窗 \(m\in\{6,8,10,12\}\) 上刚化为
\[
d_{\min}(m)=F_{m/2},
\qquad
N_{\min}(m)=F_m,
\qquad
\lim_{g\to\infty}F_{m/2}^{\,2g-2}Z_m(\Sigma_g)=F_m.
\]
另一方面，对每个 \(a>a_{\mathrm c}\)，零温冻结极限保留顶端几何：
\[
P_a=a\alpha_\ast+g_\ast,
\qquad
\lim_{m\to\infty}\frac1m H(\pi_{m,a})=g_\ast,
\qquad
\pi_{m,a}(A_m^{\max})\to 1.
\]
而 \(\infty\)-预算端点却只剩
\[
\mathcal C_m^{\mathrm{cont}}(\infty)=2^m,
\qquad
R_m(\infty)=0,
\]
从而抹去全部顶端谱几何。

因此，同一纤维直方图同时携带一对对偶层析：高 genus 极限读取最小扇区，零温冻结读取最大扇区；而 \(\infty\)-预算端点并不是零温的对偶，而是一个把两端几何一并压平的退相干端点。
\end{theorem}
\begin{proof}
最小扇区层析由定理 \ref{thm:conclusion-foldbin-highgenus-minsector-tomography} 给出，而审计偶窗上的 Fibonacci 精确值则由推论 \ref{cor:conclusion-foldbin-audited-even-fibonacci-minsector-genus-recovery} 给出。

零温冻结对最大层的指数集中、压力公式与熵率极限，正是定理 \ref{thm:conclusion-frozen-escort-shannon-renyi-collapse} 的内容。最后，\(\infty\)-预算端点只保留总质量且与零温冻结不交换，则由定理 \ref{thm:conclusion-freezing-infinity-budget-noncommutation} 直接给出。证毕。
\end{proof}

\begin{theorem}[跨尺度样本相等模式族闭合热相 Legendre 接触几何，并在开区间上局部刚性]\label{thm:conclusion-sample-pattern-hot-contact-local-rigidity}
记 \(\Pi_{|X_m|}\) 为按推前权重抽取的 \(|X_m|\) 个样本的相等模式分割。则跨尺度族
\[
\{\mathrm{Law}(\Pi_{|X_m|})\}_{m\ge 1}
\]
唯一决定热相中的全部对象：
\[
f(a),\qquad
\tau(q),\qquad
q_\ast(a),\qquad
a\longleftrightarrow q_\ast(a).
\]
更进一步，若两套系统在某开区间 \(J\subset(a_0,a_1)\) 上满足
\[
\Gamma^{\mathrm{orc}}_1(a)=\Gamma^{\mathrm{orc}}_2(a)
\qquad(a\in J),
\]
则不仅
\[
f_1(a)=f_2(a)
\qquad(a\in J),
\]
而且在相应 \(q\)-像区间上有
\[
\tau_1(q)=\tau_2(q).
\]
故热相预言机容量指数不是粗略投影，而是局部刚性的完整接触不变量。
\end{theorem}
\begin{proof}
跨尺度样本模式族对 \(\Gamma^{\mathrm{orc}}\)、\(f\)、\(\tau\) 与 \(a\leftrightarrow q_\ast(a)\) 的全恢复，正是定理 \ref{thm:conclusion-sample-pattern-family-closes-hot-contact-geometry} 的内容。

若 \(\Gamma^{\mathrm{orc}}_1=\Gamma^{\mathrm{orc}}_2\) 在开区间 \(J\) 上逐点一致，则其一、二阶导数在 \(J\) 上也一致。由同一定理中的互反公式
\[
q_\ast(a)=1-\Gamma_{\mathrm{orc}}'(a),
\qquad
\tau(q_\ast(a))=\Gamma_{\mathrm{orc}}(a)-a\Gamma_{\mathrm{orc}}'(a)
\]
可知两套系统在 \(J\) 上有相同的 \(q_\ast(a)\) 与 \(f(a)=\Gamma^{\mathrm{orc}}(a)-a\)，并在相应 \(q\)-像区间上有相同的 \(\tau(q)\)。证毕。
\end{proof}

\begin{theorem}[连续容量曲线与主序/InfoNCE 塔构成固定分辨率下的普适全息对象]\label{thm:conclusion-capacity-holography-majorization-infonce-universal-object}
固定分辨率 \(m\)。则连续容量曲线
\[
\mathcal C_m^{\mathrm{cont}}(u)\qquad(u\ge 0)
\]
与下列对象彼此完全等价：
\[
\{h_d^{(m)}\}_{d=1}^{2^m},
\qquad
N_m(T),
\qquad
S_0(m),S_1(m),\dots,S_{2^m-1}(m),
\]
\[
(w_1,\dots,w_{|X_m|}),
\qquad
\{\mathcal L_{K,m}^{\star}\}_{K=2}^{|X_m|},
\]
以及固定尺度的 Wedderburn 块型
\[
A_m\cong \bigoplus_{d\ge 1}M_d(\CC)^{\oplus n_d(m)}.
\]
特别地，在固定总质量切片上，
\[
\mathbf d\succcurlyeq_{\mathrm{HLP}}\mathbf e
\iff
\mathcal C_{\mathbf d}(u)\le \mathcal C_{\mathbf e}(u)\ \ (\forall u\ge 0)
\iff
\mathscr L_{\mathbf d}\preccurlyeq_{\mathrm{InfoNCE}}\mathscr L_{\mathbf e}.
\]
故连续容量曲线并非“预算摘要”，而是固定分辨率下同时编码直方图、尾计数、有限矩列、有序权重谱、完整 Bayes--InfoNCE 塔与半单块结构的单一普适坐标；而主序只是这一坐标在固定总质量切片上的一个等价偏序投影。
\end{theorem}
\begin{proof}
连续容量曲线对有序纤维谱、完整 Bayes--InfoNCE 数值塔与 Wedderburn 型的全息完备性，正是定理 \ref{thm:conclusion-capacity-curve-order-holography-universality} 的内容。固定总质量切片上主序、容量逐点序与完整 Bayes--InfoNCE 塔之间的严格反向同构，则由定理 \ref{thm:conclusion-majorization-capacity-infonce-strict-antiisomorphism} 给出。二者合并即得。证毕。
\end{proof}

\begin{theorem}[纤维多重度 Mellin--Dirichlet 主核是算术--拓扑双面的共同完备对象]\label{thm:conclusion-mellin-master-kernel-arithmetic-topological-complete-object}
\leanverified{paper\_conclusion\_mellin\_master\_kernel\_arithmetic\_topological\_complete\_object}
对每个固定窗口 \(m\)，定义
\[
\nu_m(d):=\#\{x:d_m(x)=d\},
\qquad
M_m(s):=\sum_{d\ge 1}\nu_m(d)d^{-s}.
\]
则
\[
M_m(-q)=S_q(m)\qquad(q\ge 1),
\]
\[
M_m(2g-2)=Z_m(\Sigma_g)\qquad(g\ge 1).
\]
更进一步，任意一侧的无穷采样已经唯一决定直方图 \(\nu_m\)。再把窗口方向加入，直方图塔 \((\nu_m)_m\) 同时唯一决定：
\[
\mathbf C[R_m^{\mathrm{bin}}]\ \text{的 Wedderburn 分解},
\qquad
Z_m(\Sigma_g)\ \text{对全部 }g,
\]
\[
|G_m^{\mathrm{bin}}|=\prod_d(d!)^{\nu_m(d)},
\qquad
(C_m)_m,\qquad
(B_{2r})_{r\ge 1},\qquad
(\zeta(2r))_{r\ge 1}.
\]
因此，\(M_m(s)\) 不是碰撞矩与 genus 振幅的两张并列表，而是唯一的主核：其负侧读取统计热量，其正侧读取拓扑振幅，而跨窗口塔则把半单代数、规范体积与 even-\(\zeta\) 算术同步锁死。
\end{theorem}
\begin{proof}
这正是定理 \ref{thm:conclusion-mellin-master-kernel-two-sided-completeness} 的内容。其证明已经把双侧整数点采样、有限支撑 Dirichlet 多项式的一侧刚性，以及直方图塔对 Wedderburn 块谱、规范群规模与 Bernoulli--\(\zeta\) 校正链的完备控制一次性闭合。证毕。
\end{proof}

\begin{theorem}[热力学广延层的完全扁平与 even-\texorpdfstring{$\zeta$}{zeta} 算术塔的次广延滞留]\label{thm:conclusion-thermodynamic-flatness-evenzeta-subextensive-reservoir}
设 \(\mu_m\) 为 bin-fold 推前测度。则对任意 \(q>0\)（\(q=1\) 按 Shannon 延拓解释）都有
\[
\lim_{m\to\infty}\frac1m H_q(\mu_m)=\log\varphi.
\]
换言之，全部 R\'enyi 热力学在广延斜率层面塌缩到同一常数 \(\log\varphi\)。

然而规范重整化序列却保留完整的 Stirling--Bernoulli 渐近，因而每个 \(B_{2r}\) 与 \(\zeta(2r)\) 都可由 \(\log C_m\) 的次广延校正逐阶恢复。更强地，若两套系统在每个 \(m\) 上连续容量曲线相同，则不仅全部 \(S_q(m)\)、\(H_q(\mu_m)\) 一致，而且整条
\[
(C_m),\qquad
(B_{2r}),\qquad
(\zeta(2r))
\]
校正塔也逐项一致。

因此，该体系存在严格的二层分裂：广延热力学只有单一斜率 \(\log\varphi\)，而完整的 even-\(\zeta\) 算术信息则整体沉积在次广延规范层，并且被容量等价类 rigidify。
\end{theorem}
\begin{proof}
广延 R\'enyi 轴全部塌缩为 \(\log\varphi\)，而 even-\(\zeta\) 算术塔整体滞留于次广延层，正是定理 \ref{thm:conclusion-thermodynamic-flattening-subextensive-arithmetic-rigidity} 的内容。该定理还同时说明：一旦连续容量曲线逐窗相同，则广延热力学层与 Bernoulli--\(\zeta\) 校正层都会被同步 rigidify。证毕。
\end{proof}

\begin{theorem}[余数剖面的模 \texorpdfstring{$3$}{3} 二分与第一坐标偏置的 Fibonacci 共振忠实性]\label{thm:conclusion-residue-parity-blindmode-localbias-fibonacci-resonance}
设
\[
M=F_{m+2},
\qquad
d_m(r)=\#\{\omega:r(\omega)=r\}.
\]
若 \(M\) 为奇数，等价于
\[
m\not\equiv 1\pmod 3,
\]
则 \(d_m(\cdot)\) 唯一决定一切有限坐标模式计数
\[
A_{m,I}(r;v),
\]
特别地唯一决定全部单坐标条件概率 \(p_k(r)\) 及完整推前核 \(P_m\)。若
\[
m\equiv 1\pmod 3,
\]
则最简单切片 \(k=1\) 的逆问题恰有一维核
\[
\ker(I+T_1)=\CC\chi_{M/2},
\qquad
\chi_{M/2}(r)=(-1)^r,
\]
故偶窗口中的最小局部不可逆性不是高维噪声，而是严格的一维 parity 盲模。

另一方面，定义第一坐标偏置能量
\[
\nu_{m,1}(r):=\mu_m(r)\bigl(2p_1(r)-1\bigr),
\qquad
E_{m,1}:=\sum_r \nu_{m,1}(r)^2.
\]
则有严格等价
\[
E_{m,1}=0
\iff
d_m(r)\equiv \frac{2^m}{F_{m+2}}
\iff
\mathrm{Col}_m=\frac1{F_{m+2}}.
\]
并且
\[
\liminf_{m\to\infty}F_{m+2}E_{m,1}
\ge
\frac12\tan^2(\pi\varphi^{-2})\,C_\varphi^2>0.
\]
驱动该下界的主导频率与全局碰撞缺口所用者完全相同：
\[
k=F_m,\qquad k=F_{m+1}.
\]
因此，第一坐标偏置不是全局缺口的旁证，而是同一 bulk Fibonacci 共振在局部一维投影下的忠实显形。
\end{theorem}
\begin{proof}
奇窗口全反演与偶窗口一维 parity 盲模的精确二分，正是定理 \ref{thm:conclusion-oddwindow-full-inversion-evenwindow-rankone-parity-blindmode} 的内容。第一坐标偏置能量对全局均匀性的忠实性，以及它与碰撞缺口共用同一对 Fibonacci 共振频率，则由定理 \ref{thm:conclusion-first-coordinate-bias-global-collision-resonance-homology} 给出。证毕。
\end{proof}

\begin{conjecture}[固定尺度 bin-fold 的最小完备审计对象是“容量曲线 + parity 补账”二层对]\label{conj:conclusion-binfold-minimal-complete-object-capacity-plus-parity-ledger}
对每个固定分辨率 \(m\)，存在有限维奇偶角色空间 \(B_m\)，满足
\[
B_m=0\iff m\not\equiv 1\pmod 3,
\]
而当 \(m\equiv 1\pmod 3\) 时，\(B_m\) 恰由全部局部逆问题的盲模角色生成。并且一旦附加一份补账
\[
\beta_m\in B_m^\vee,
\]
则二元数据
\[
\bigl(\mathcal C_m^{\mathrm{cont}},\beta_m\bigr)
\]
足以唯一恢复全部有限坐标模式计数 \(A_{m,I}(r;v)\) 与完整推前核 \(P_m\)。

该猜想的含义是：固定尺度的最小完备对象不应是单独的容量曲线，而应是“Schur/统计全息层 + parity/局部补账层”的二层审计对。
\end{conjecture}
\begin{proof}
定理 \ref{thm:conclusion-capacity-holography-majorization-infonce-universal-object} 已说明：连续容量曲线 \(\mathcal C_m^{\mathrm{cont}}\) 唯一控制全部全局可见统计、容量序与半单块谱。另一方面，定理 \ref{thm:conclusion-residue-parity-blindmode-localbias-fibonacci-resonance} 表明：当 \(m\equiv 1\pmod 3\) 时，偶窗口中的最小局部不可逆性严格压缩为一维 parity 盲模；而在奇窗口中该补账空间消失。故把“全局容量坐标”与“局部 parity 补账”组合成最小完备对象，是当前稿内最自然的终端闭包形式。证毕。
\end{proof}

\begin{theorem}[有限缺陷剖面的原子数等于其内禀最小 Cauchy 核秩]\label{thm:conclusion-finite-defect-atomic-count-equals-intrinsic-minimal-cauchy-rank}
\leanverified{paper\_conclusion\_finite\_defect\_atomic\_count\_equals\_intrinsic\_minimal\_cauchy\_rank}
设有限缺陷剖面
\[
H(x)=\sum_{j=1}^{\kappa}\frac{w_j}{(x-\gamma_j)^2+a_j^2},
\qquad
w_j>0,\ a_j>0,
\]
且 \((\gamma_j,a_j)\) 两两不同。若存在某个 Cauchy 型正核
\[
K(x,y)=\sum_{\ell=1}^{r}
\frac{\omega_\ell}{(x-\alpha_\ell+\mathrm{i}b_\ell)(y-\alpha_\ell-\mathrm{i}b_\ell)}
\qquad
(\omega_\ell>0,\ b_\ell>0)
\]
满足
\[
K(x,x)=H(x)
\qquad(\forall x\in\RR),
\]
则必有
\[
r\ge \kappa.
\]
反之，恰存在一个 \(\kappa\)-项 Cauchy 正核精确实现该对角剖面。故 \(\kappa\) 不是拟合表示中的坐标选择产物，而是边界剖面 \(H\) 自身强制给出的内禀最小核秩。
\end{theorem}
\begin{proof}
这正是定理 \ref{thm:conclusion-finite-defect-profile-minimal-cauchy-rank} 与推论 \ref{cor:conclusion-finite-defect-profile-intrinsic-minimal-rank} 的内容。前者证明任何精确核化都必须至少使用 \(\kappa\) 个 Cauchy 原子，后者则把这一最小秩解释为边界剖面自身的内禀不变量。证毕。
\end{proof}

\endinput
